
One of the major challenge of this research was identifying net and gross capital stock measures of the non-residential sector for Chile, this is the first contribution to provide both types of measures under a stock-flow consistent framework. To do so, it was necessary to leverage on different sources, mainly the seminal contributions of \cite{Hofman2000} that are complemented with two interrelated research agendas about Chile's Economic History. On the one hand, the works of Clio-Lab PUC historical statistics series presented in \citet{Luders2016}, which to date, are the data more comprehensive publicly available historical statistics for Chile. It should be noticed that this reference is the presentation of multiple working papers on Chilean historical statistics, where the methodological developments on capital stocks and investments can be found in \citet{Diaz2016}. This measurements are all net measures capital stock without identifying the non-residential construction component of infrastucture capital stocks in distinction to residential ones. 

On the other hand, my data base also leverages on contributions about capital stock measures developed by the Economic History school of the Barcelona University, mainly in the works on gross capital stock series in Machinery and Equipment and Non-Residential Construction sector set in motion by the research program on Latin America modernization \citep{Carreras2006}. The most contemporanous advance from this body of research is synthesized by \citet{Tafunell2016} who provide a gross capital stocks series for the non-residential sector and for machinery and equipment. 

In this appendix is presented the procedure on how the aforementioned references have been systematized and articulated, to rely in new series of capital stocks with net and gross measures. Following \cite{BasuRama2013} comparisons between both types of series are useful because they provide robustness to empirical findings of profitability. This is consistent with \citet{Shaikh2016} claims that neoclassical theory is prone to use net measures instead of gross because their interest only on valuation, while a Marxist interested in the dynamics of capital accumulation understood as the actual levels of capital stocks instead of their book value, would prefer to look at gross measures. 

\cite{Shaikh2016} provides an interesting discussion about the quality adjustments introduced in fixed asset measures provided by the US Bureau of Economic Analysis, arguing that chain weighted measurements that have become an standard in OECD countries are problematic because they break the stock-flow consistency between investments and capital stocks having as a consequence distortions in their use when doing a transition from real to nominal measures. Facing this challenge \cite{Shaikh2016} introduces a Generalized Perpetual Inventory Method (GPIM) that allows him to transit from chain weighted measurements of real capital stocks, to an adjustment of real capital stocks that are not distorted by quality adjustments. Furthremore, Shaikh's GPIM allows to transit between chain weighted, to historical, and replacement costs measurements in real and nominal terms. In my assesment, this is not so different to \citet{Basu2013} proof of when does it matters to use historical or replacement cost measurements of capital stocks: only when there is a rate of change of capital good prices over a certain threshold. Otherwise, there is no relevance between them. This have been also discusssed more recently by \citet{Marquetti2010} in a context that actually matters: the case of Brazil which faced periods of severe inflation in capital goods throughout during the import substitution industrialization strategy period. 

Despite that Shaikh have consistently vouched for a Standard Interpretion (SI) of the transformation problem of value to prices, my assessment is that his GPIM is also useful for authors as myself who locate their research within the New Interpretation (NI) of the transformation problem. This because the NI entails a theory of money through the monetary expression of labour value, considering that value express itself in money capital throughout the capital circuit. Hence, to sustain stock-flow consistency when working with real measures is necessary to transit to nominal measures which are the appropiate measures of profitability. Furtheremore, the GPIM also provides a connection between gross and net measures, which is important for the case of sustaining stock-flow consistency of my own contribution with this data base.  While \citet{Shaikh2016} discusses a more elaborated framework of the transition between chain weighted real capital stock measures to unadjusted ones, here I present a more simple and intuitive approach that allows to check the stock-flow consistency of my data as follows:

Let the equation of motion of gross capital stocks in real terms be: 

\begin{align}
\Delta Kg_{t} &= Ig_{t} -R_{t}\  
\end{align}

Where $\Delta Kg_{t}$ refers to the accumulation of gross capital stock, $Ig_{t}$ to gross fixed capital formation or gross investment flow, and $R_{t}$ to retirements of capital stocks. As pointed out for the case of Chile (for a more contemporaneous period that the one of my interest, 1986-2015) by \citet{Ffrench2016} in a similar fashion of \citet{Shaikh2016} concerns, capital retirements are a  relevant concept because identify the movement on the actual installed capacity instead of book valuation as net measures does. 

Accordingly, let the equation of motion of net capital stocks in real terms be:  

\begin{align}
\Delta Kn_{t} &= Ig_{t} -D_{t}    
\end{align}

Where $Kn_{t}$ is the net measure of capital stock. On the other hand, $D_{t}$ is consumption of fixed capital in each period which is equivalent to: 

\begin{align}
D_{t} &= \delta_{t} Kn_{t}
\end{align}


Where $\delta_{t}$ is the annual rate of depreciation of net capital stocks. While this is not intuitive in terms of the life of different types of tangible assets, its level, average, standard deviation, and trend does delivers information to check if of leveraging on different estimates of gross and net capital stock measures are plausible estimations under a stock-flow consistent rule. 

With both equations of motion for gross and net capital stocks equated to $Ig_{t}$ is straightforward to identify: $\delta_{t}$ as: 

\begin{align}
\delta_{t} = \dfrac{\Delta Kg_{t} - \Delta Kn_{t} + R_{t}}{Kn_{t}}    
\end{align}

In the following table some summary statistics of the implicit depreciation rates are presented, their results are discussed in the following subsections, for each one of the measurements of gross and net capital stocks by types: machinery and equipment ($ME$), construction sector ($C$), non-residential construction sector ($NRC$), residential construction sector ($RC$), total ($T$), and non-residential capital stocks ($NR$). Time series plots for implicit annual depreciation rates are also presented in each section. 

\input{AppendixA_Tables/AppendixA_Kstock_appendix_Z_rates}

Correspondingly, defining the depletion rate of capital retirements in relation to gross stocks as: 

\begin{align}
R_{t} &= z_{t} Kn_{t}
\end{align}


It is also straightforward to identify the depletion rate of capital retirements as: 

\begin{align}
z_{t} = \dfrac{\Delta Kn_{t} - \Delta Kg_{t} + D_{t}}{Kg_{t}}
\end{align}


In the following table some summary statistics of the implicit depreciation rates are presented, their results are discussed in the following subsections, for each one of the measurements of gross and net capital stocks by types: machinery and equipment ($ME$), construction sector ($C$), non-residential construction sector ($NRC$), residential construction sector ($RC$), total ($T$), and non-residential capital stocks ($NR$). Time series plots for implicit annual depletion rates of gross capital stocks are also presented in each section. 

\input{AppendixA_Tables/AppendixA_Kstock_appendix_Z_rates}

In the following subsections the sources of different measures of capital stocks and the procedure under this simple stock-flow consistent rule for the estimation of each component is presented. It should be noticed that the annual depreciation rate does not shows information on the life of tangible assets as standard PIM rules leverage for estimations on capital stock measures. Instead, it collapses the dynamics of net tangible asssets accumulation and replacements that is implicit by leveraging in multiple sources.  

\section{Measurements of Capital Stocks in Machinery and Equipment}

For the case of capital stock measurements in machinery and equipment ${K^{g}}_{ME}$, Perez (2019) extension of Lüders et al. (2016) have included gross fixed capital formation series available in Tafunell et al. (2015) in its exended version of chilean historical statistics. Hence, to sustain stock flow consistency for net and gross capital stock measurements of machinery and equipment I leverage on Tafunell index for real gross capital stock in machinery and equipment, using only a point estimate of Hofman (2000) for the year 1950 to turn the index into levels.  An important detail worth to mention is that Hofman series are in 1980 chilean pesos, while Perez (2019) are in 2003 millions of pesos. Hence, before estimating gross capital stock using Hofman's point estimate for the year 1950 and Tafunell index, the 1950 point estimate is deflated to 2003 thousands of millions of pesos (instead of millions as Perez only for more parsimonious visualization of the data while working with it). 
	
Utilizes gross fixed capital formation (${I^{g}}_{ME}$=), data from the consolidated dataset by Perez (2019) and Clio-Lab. This data is harmonized with Tafunell \& Ducoing (2015) indices to estimate consistent gross fixed capital formation aggregates aligned with sectoral GDP estimates for the 1860-1940 period, thereby extending Chile's national accounts historical statistics. Retirements of machinery and equipment ($R_{ME}$), are determined by the formula presented previously $R_{ME} = Ig_{ME} - \Delta Kg_{ME}$. While the implicit depletion rate of machinery and equipment ($z_{ME}$) are determined by the stock-flow consistency formula, and its estimations can be appreciated in the figure below. 

\begin{figure}[H]
    \centering
    \caption{\centering Implicit Annual Depletion Rate of Gross Capital Stocks, \\ Machinery and Equipment (1901–1994).}
    \includegraphics[width=0.5\textwidth]{AppendixA_Figures/AppendixA_z_ME_Kstock_appendix_plot_Zrates.png}
    \begin{minipage}{0.5\textwidth}
        \centering
        \footnotesize
        \tikz[baseline]{\draw[darkgreen, thick] (0,0)--(0.5,0);}~Linear trend (dark green) \hspace{1em}
        \tikz[baseline]{\draw[blue, thick] (0,0)--(0.5,0);}~LOESS trend (blue)
    \end{minipage}
    \label{fig:d_ME_Kstock_appendix_plot_A}
\end{figure}


As \citet{Tafunell2016} provide only an index of gross capital stock, net capital stock (${K^{n}}_{ME}$) calculations are sourced from Clio-Lab. \citet{PerezEyzaguirre2019} relies on gross fixed capital formation for machinery and equipment without discrepancies when incorporating net capital stocks estimated via a perpetual inventory method, as presented by \citet{Diaz2016}. However, to sustain level consistency levels, I rely on \citet{Hofman2000} point estimate of net capital stocks in $ME$ and used Clio-Lab PUC data to build an index of real net capital stocks at year base $= 100$ to construct the series. Total fixed capital consumption ($D_{ME}$) is identified by the difference between gross fixed capital formation in machinery and equipment with the annual variation of net capital stocks, i.e $D_{t} = {K^{g}}_{t}- {K^{n}}_{t}$. This series is necessary to be estimated for estimating net output and net profit income measures. 

The implicit annual depreciation rate is calculated with the formula presented previously in the introduction of the appendix. The average rate registers in -1.6\%, a standard deviation of 3.9\%, and the slope of its time trend registers a negative slope. This results are intuitive once considering that negative values in average might dominate to positive ones, given that in some periods the depreciation of machinery might rise or decrease the value of net assets in bumpy and asymmetric patterns. Moreover, the negative slope of its time trend, is consistent with the fact that technological progress foster longer duration of machinery and equipment, but also higher pace of their replacements with technological innovations. A time series plot for each implicitly annual depreciation rate of net stocks in machinery and equipment is presented in the figure below of this appendix subsection. 

\begin{figure}[H]
    \centering
    \caption{\centering Implicit Annual Depreciation Rate of Net Capital Stocks, \\ Machinery and Equipment (1901–1994).}
    \includegraphics[width=0.5\textwidth]{AppendixA_Figures/AppendixA_d_ME_Kstock_appendix_plot_Drates.png}
    \begin{minipage}{0.5\textwidth}
        \centering
        \footnotesize
        \tikz[baseline]{\draw[darkgreen, thick] (0,0)--(0.5,0);}~Linear trend (dark green) \hspace{1em}
        \tikz[baseline]{\draw[blue, thick] (0,0)--(0.5,0);}~LOESS trend (blue)
    \end{minipage}
    \label{fig:d_ME_Kstock_appendix_plot_B}
\end{figure}

%For the case of capital stock measurements in machinery and equipment ${K^{g}}_{ME}$, Perez (2019) extension of Lüders et al. (2016) have included gross fixed capital formation series available in Tafunell et al. (2015) in its exended version of chilean historical statistics. Hence, to sustain stock flow consistency for net and gross capital stock measurements of machinery and equipment I leverage on Tafunell index for real gross capital stock in machinery and equipment, using only a point estimate of Hofman (2000) for the year 1950 to turn the index into levels.  An important detail worth to mention is that Hofman series are in 1980 chilean pesos, while Perez (2019) are in 2003 millions of pesos. Hence, before estimating gross capital stock using Hofman's point estimate for the year 1950 and Tafunell index, the 1950 point estimate is deflated to 2003 thousands of millions of pesos (instead of millions as Perez only for more parsimonious visualization of the data while working with it). 
	
Utilizes gross fixed capital formation (${I^{g}}_{ME}$=), data from the consolidated dataset by Perez (2019) and Clio-Lab. This data is harmonized with Tafunell \& Ducoing (2015) indices to estimate consistent gross fixed capital formation aggregates aligned with sectoral GDP estimates for the 1860-1940 period, thereby extending Chile's national accounts historical statistics. Retirements of machinery and equipment ($R_{ME}$), are determined by the formula presented previously $R_{ME} = Ig_{ME} - \Delta Kg_{ME}$. While the implicit depletion rate of machinery and equipment ($z_{ME}$) are determined by the stock-flow consistency formula, and its estimations can be appreciated in the figure below. 

\begin{figure}[H]
	\centering
	\caption{Implicit Annual Depletion Rate of Gross Capital Stocks. Machinery and Equipment (1901-1994).}
	\begin{subfigure}{0.5\textwidth}
		\centering
		\includegraphics[scale=0.5]{AppendixA_Figures/AppendixA_z_ME_Kstock_appendix_plot_Zrates.png}
	\end{subfigure}
	\hfill
	\caption*{\tikz\draw[darkgreen, thick] (0,0) -- (0.5,0) node[right] {(dark green)}; Linear Time Trend}
	\hfill
	\caption*{\tikz\draw[blue, thick] (0,0) -- (0.5,0) node[right] {(blue)}; LOESS function Time Trend}
	\label{fig: d_ME_Kstock_appendix_plot_A}
\end{figure}

As \citet{Tafunell2016} provide only an index of gross capital stock, net capital stock (${K^{n}}_{ME}$) calculations are sourced from Clio-Lab. \citet{PerezEyzaguirre2019} relies on gross fixed capital formation for machinery and equipment without discrepancies when incorporating net capital stocks estimated via a perpetual inventory method, as presented by \citet{Diaz2016}. However, to sustain level consistency levels, I rely on \citet{Hofman2000} point estimate of net capital stocks in $ME$ and used Clio-Lab PUC data to build an index of real net capital stocks at year base $= 100$ to construct the series. Total fixed capital consumption ($D_{ME}$) is identified by the difference between gross fixed capital formation in machinery and equipment with the annual variation of net capital stocks, i.e $D_{t} = {K^{g}}_{t}- {K^{n}}_{t}$. This series is necessary to be estimated for estimating net output and net profit income measures. 

The implicit annual depreciation rate is calculated with the formula presented previously in the introduction of the appendix. The average rate registers in -1.6\%, a standard deviation of 3.9\%, and the slope of its time trend registers a negative slope. This results are intuitive once considering that negative values in average might dominate to positive ones, given that in some periods the depreciation of machinery might rise or decrease the value of net assets in bumpy and asymmetric patterns. Moreover, the negative slope of its time trend, is consistent with the fact that technological progress foster longer duration of machinery and equipment, but also higher pace of their replacements with technological innovations. A time series plot for each implicitly annual depreciation rate of net stocks in machinery and equipment is presented in the figure below of this appendix subsection. 

\begin{figure}[H]
\centering
\caption{Implicit Annual Depreciation Rate of Net Capital Stocks. Machinery and Equipment (1901-1994).}
\begin{subfigure}{0.5\textwidth}
	\centering
	\includegraphics[scale=0.5]{AppendixA_Figures/AppendixA_d_ME_Kstock_appendix_plot_Drates.png}
\end{subfigure}
\hfill
\caption*{\tikz\draw[darkgreen, thick] (0,0) -- (0.5,0) node[right] {(dark green)}; Linear Time Trend}
\hfill
\caption*{\tikz\draw[blue, thick] (0,0) -- (0.5,0) node[right] {(blue)}; LOESS function Time Trend}
\label{fig: d_ME_Kstock_appendix_plot_B}
\end{figure}

\section{Measurements of Capital Stocks in Construction Sector}

In the case of the construction sector as a whole (${K^{g}}_{C}$), Tafunell \& Ducoing (2015) do not provide an index on the construction sector as a whole, rather only on the Non-Residential sector. An assesment on the series based on GFCF index by Tafunell (2013) points out that the series is build leveraging on data of cement inputs\footnote{I'm grateful to Ignacio Perez for pointing this out to me, as every discussion on proper measures to sustain stock-flow consistency with his data set.}. While this is not an appropriate for a measure of the construction sector as a whole because does not consider more basic materials of residential construction, as for example adobe houses, it is plausible  that the non-residential sector was intensive on cement inputs. Hence, it is plausible to assume that their estimates are appropriate to measure the non-residential construction sector.  Nevertheless they are not for the construction sector as a whole. Hence, my approach was to sustaint the consistency between gross and net capital stocks by means of identifying fixed capital consumption in the construction sector implicit in the net stock series of Diaz \& Wagner (2016) presented in the Clio-Lab PUC series and its corresponding extension spliced by Perez (2019). To do so, I rely on the accounting equations of gross and net capital stocks previously presented, identifying the total fund of fixed capital consumption to later identify the gross capital stock for the construction sector as a whole. Having point estimates on gross capital stocks and consumption of fixed capital of the construction sector, while complete series of fixed capital formation and net capital stocks, then the equations of motions for both measures and the presented identities at the begining of this appendix, identified allows me to calculate iteratively total fixed capital consumption and gross capital with the cost of losing the initial period $t=1900$.
	
The gross fixed capital formation data in the construction sector (${I^{g}}_{C}$) from the consolidated data set by Perez (2019) and Clio-Lab are available. In this case, he does not rely on Tafunell \& Ducoing (2015) indexes to estimate gross fixed capital formation in the construction sector due to the aforementioned assessment. Therefore, I keep following stock flow consistency to build gross and net capital stocks with Perez (2019) splice with Lüders et al. (2016) data. Consequently, retirements of machinery and equipment ($R_{C}$) are determined by the formula $R_{C} = Ig_{C} - \Delta {K^{g}}_{C}$, ex-post estimating gross capital stocks. Implicit annual depletion rate of the construction sector determined by the formula presented above ($z_{C}$), as can be appreciated in the figure below. 

\begin{figure}[H]
    \centering
    \caption{\centering Implicit Annual Depletion Rate of Gross Capital Stocks, \\ Machinery and Equipment (1901–1994).}
    \includegraphics[width=0.5\textwidth]{AppendixA_Figures/AppendixA_z_ME_Kstock_appendix_plot_Zrates.png}
    \begin{minipage}{0.5\textwidth}
        \centering
        \footnotesize
        \tikz[baseline]{\draw[darkgreen, thick] (0,0)--(0.5,0);}~Linear trend (dark green) \hspace{1em}
        \tikz[baseline]{\draw[blue, thick] (0,0)--(0.5,0);}~LOESS trend (blue)
    \end{minipage}
    \label{fig:d_ME_Kstock_appendix_plot_A}
\end{figure}

For the case of net capital stocks of the construction sector (${K^{n}}_{C}$), I rely on the series of capital stock presented in Clio-Lab PUC historical statistics (Lüders, et al., 2016), estimated under a perpetual inventory method by \citet{Diaz2016}. As previosuly mentioned, the total fixed capital consumption for the construction sector ($D_{C}$) is identified by the difference between gross fixed capital formation in construction sector with the annual variation of net capital stocks, i.e $D_{t} = {I^{g}}_{t}- \Delta {K^{n}}_{t}$. Correspondingly, the implicit annual depreciation rate ($d_{C}$) is calculated with the identity presented in the introduction of the appendix, with an average of 2.47\% and standard deviation of 0.73\%, while presenting a positive slope in its linear time trend at a value of 1.17 as can be appreciated in the table ~\ref{table:depraciation_rates}, presented above at the introduction of the appendix. 

\begin{figure}[H]
    \centering
    \caption{\centering Implicit Annual Depreciation Rate of Net Capital Stocks, \\ Total Construction Sector (1901–1994).}
    \includegraphics[width=0.5\textwidth]{AppendixA_Figures/AppendixA_d_C_Kstock_appendix_plot_Drates.png}
    \begin{minipage}{0.5\textwidth}
        \centering
        \footnotesize
        \tikz[baseline]{\draw[darkgreen, thick] (0,0)--(0.5,0);}~Linear trend (dark green) \hspace{1em}
        \tikz[baseline]{\draw[blue, thick] (0,0)--(0.5,0);}~LOESS trend (blue)
    \end{minipage}
    \label{fig:d_C_Kstock_appendix_plot}
\end{figure}


%In the case of the construction sector as a whole (${K^{g}}_{C}$), Tafunell \& Ducoing (2015) do not provide an index on the construction sector as a whole, rather only on the Non-Residential sector. An assessment on the series based on GFCF index by Tafunell (2013) points out that the series is build leveraging on data of cement inputs\footnote{I'm grateful to Ignacio Perez for pointing this out to me, as every discussion on proper measures to sustain stock-flow consistency with his data set.}. While this is not an appropriate for a measure of the construction sector as a whole because does not consider more basic materials of residential construction, as for example adobe houses, it is plausible  that the non-residential sector was intensive on cement inputs. Hence, it is plausible to assume that their estimates are appropriate to measure the non-residential construction sector.  Nevertheless they are not for the construction sector as a whole. Hence, my approach was to sustaint the consistency between gross and net capital stocks by means of identifying fixed capital consumption in the construction sector implicit in the net stock series of Diaz \& Wagner (2016) presented in the Clio-Lab PUC series and its corresponding extension spliced by Perez (2019). To do so, I rely on the accounting equations of gross and net capital stocks previously presented, identifying the total fund of fixed capital consumption to later identify the gross capital stock for the construction sector as a whole. Having point estimates on gross capital stocks and consumption of fixed capital of the construction sector, while complete series of fixed capital formation and net capital stocks, then the equations of motions for both measures and the presented identities at the begining of this appendix, identified allows me to calculate iteratively total fixed capital consumption and gross capital with the cost of losing the initial period $t=1900$.
	
The gross fixed capital formation data in the construction sector (${I^{g}}_{C}$) from the consolidated data set by Perez (2019) and Clio-Lab are available. In this case, he does not rely on Tafunell \& Ducoing (2015) indexes to estimate gross fixed capital formation in the construction sector due to the aforementioned assessment. Therefore, I keep following stock flow consistency to build gross and net capital stocks with Perez (2019) splice with Lüders et al. (2016) data. Consequently, retirements of machinery and equipment ($R_{C}$) are determined by the formula $R_{C} = Ig_{C} - \Delta {K^{g}}_{C}$, ex-post estimating gross capital stocks. Implicit annual depletion rate of the construction sector determined by the formula presented above ($z_{C}$). As can be appreciated in the figure below. 


\begin{figure}[H]
	\centering
	\caption{Implicit Annual Depletion Rate of Gross Capital Stocks. Total Construction sector (1901-1994).}
	\begin{subfigure}{0.5\textwidth}
		\centering
		\includegraphics[scale=0.5]{AppendixA_Figures/AppendixA_z_C_Kstock_appendix_plot_Zrates.png}
	\end{subfigure}
	\hfill
	\subcaption*{\tikz\draw[darkgreen, thick] (0,0) -- (0.5,0) node[right] {(dark green)}; Linear Time Trend}
	\hfill
	\subcaption*{\tikz\draw[blue, thick] (0,0) -- (0.5,0) node[right] {(blue)}; LOESS function Time Trend}
	\label{fig: d_ME_Kstock_appendix_plot}
\end{figure}

For the case of net capital stocks of the construction sector (${K^{n}}_{C}$), I rely on the series of capital stock presented in Clio-Lab PUC historical statistics (Lüders, et al., 2016), estimated under a perpetual inventory method by \citet{Diaz2016}. As previosuly mentioned, the total fixed capital consumption for the construction sector ($D_{C}$) is identified by the difference between gross fixed capital formation in construction sector with the annual variation of net capital stocks, i.e $D_{t} = {I^{g}}_{t}- \Delta {K^{n}}_{t}$. Correspondingly, the implicit annual depreciation rate ($d_{C}$) is calculated with the identity presented in the introduction of the appendix, with an average of 2.47\% and standard deviation of 0.73\%, while presenting a positive slope in its linear time trend at a value of 1.17 as can be appreciated in the table ~\ref{table:depraciation_rates}, presented above at the introduction of the appendix. 


\begin{figure}[H]
	\centering
	\caption{Implicit Annual Depreciation Rate of Net Capital Stocks. Total Construction sector (1901-1994).}
	\begin{subfigure}{0.5\textwidth}
		\centering
		\includegraphics[scale=0.5]{AppendixA_Figures/AppendixA_d_C_Kstock_appendix_plot_Drates.png}
	\end{subfigure}
	\hfill
	\subcaption*{\tikz\draw[darkgreen, thick] (0,0) -- (0.5,0) node[right] {(dark green)}; Linear Time Trend}
	\hfill
	\subcaption*{\tikz\draw[blue, thick] (0,0) -- (0.5,0) node[right] {(blue)}; LOESS function Time Trend}
	\label{fig: d_ME_Kstock_appendix_plot}
\end{figure}

\section{Measurements of Capital Stocks in Non Residential Construction Sector}

In the case of the non-residential construction sector (${K^{g}}_{NRC}$), \citet{Tafunell2016} does provide an index for its rate of growth in real terms, since 1900 year onwards. As previously mentioned,\citet{Tafunell2016} based on \citet{Tafunell2013} series and Ducoing (2012) doctoral dissertation, is built leveraging on data of cement inputs. While I asume this is not adequate for the construction sector as a whole, I do consider it a plausible assumption for the non-residential one. Thus, to sustain stock-flow consistency I estimate a new index for the growth of non-residential construction by using Tafunell et al. index as a proportion of the rate of growth of the gross fixed capital stock explained above, which allows also to estimate the residential sector as a reminder of the compound index implicit once considering Tafunell and Cilo-Lab PUC historical statistics. This index register an average and standard deviation statistical discrepancy with Tafunell et al (2015) of 1\% and 1.9\% respectively. To leverage on initial point estimates at $t=1900$, I assume that ${K^{n}}_{NRC}= \alpha {K^{n}}_{C}$  where $\alpha$  is the share of gross fixed capital formation in non-residential construction sector over the sum of residential and non-residential gross fixed capital formation available in Hofman (2000) table D, and follow the same assumption for the initial point estimate of total fixed capital consumption $D_{NRC} = \alpha D_{C}$.  From here, I got an initial point estimate at $t=1900$ for ${K^{g}}_{NRC}$ and then I use the compounded index constructed to sustain stock-flow consistency.

For the data of gross fixed capital investment for the non-residential sector ${I^{g}}_{NRC}$ I leverage on gross fixed capital formation data in the construction sector as a whole from the consolidated date set by \citet{PerezEyzaguirre2019}, but preponderated by the share of investments implicit in Hofman (2000) series of gross fixed capital formation corresponding to the non-residential sector. In this case, the series do not rely on \citet{Tafunell2016} indexes to estimate gross fixed capital formation in the construction sector due the aforementioned assessment above. Therefore, I keep following a stock-flow consistency to build gross and net capital stocks with \citet{PerezEyzaguirre2019} extension of Clio-Lab Chilean Historical Statistics.

Retirements of capital stocks in the non-residential sector ($R_{NRC}$) that refers to fully depreciated value of plant infrastructure is estimated with stock-flow consistency by the equation $R_{NRC} = Ig_{NRC} - \Delta {K^{g}}_{NRC}$. Implicit annual depletion rate of non-residential construction ($z_{NRC}$) is determined with stock-flow consistency, and can be appreciated in the figure below. 

\begin{figure}[H]
    \centering
    \caption{\centering Implicit Annual Depletion Rate of Gross Capital Stocks, \\ Non-Residential Construction Sector (1901–1994).}
    \includegraphics[width=0.5\textwidth]{AppendixA_Figures/AppendixA_z_NRC_Kstock_appendix_plot_Zrates.png}
    \begin{minipage}{0.5\textwidth}
        \centering
        \footnotesize
        \tikz[baseline]{\draw[darkgreen, thick] (0,0)--(0.5,0);}~Linear trend (dark green) \hspace{1em}
        \tikz[baseline]{\draw[blue, thick] (0,0)--(0.5,0);}~LOESS trend (blue)
    \end{minipage}
    \label{fig:d_ME_Kstock_appendix_plot_A}
\end{figure}

Departing from the initial point estimates for  ${K^{n}}_{NRC}$  ${D}_{NRC}$  with the assumptions presented above, and already having a series of gross capital stocks and retirements, I applied the identity coming out from the equations of motions previously mentioned and iteratively estimates of ${K^{n}}_{NRC}$ and also $D_{NRC}$. The implicit yearly depreciation rate ($d_{NRC}$) is estimated with the equation presented in the introduction of the appendix having as a result an average value of 2.16\%, a standard deviation of 0.73\%, and linear time trend slope of 0.86. This result seems consistent with the notion that depreciation of net NRC capital stocks should be characterized by higher variance in its flows of capital consumption than the construction sector as a whole, but lower than capital stocks in machinery and equipment. The figure presented below plots the time series of the depreciatation rate of net capital stocks for the non residential construction sector, suggesting an upward tendency that can be interpeted as an upward tendency of depreciation of the net value of plant infrastuctures, which is consistent with \cite{Harvey2018} understanding of built environments for the production process under a capital circuit framework. 

\begin{figure}[H]
	\centering
	\caption{Implicit Annual Depreciation Rate of Net Capital Stocks. Non-Residential Construction Sector (1901-1994).}
	\begin{subfigure}{0.5\textwidth}
		\centering
		\includegraphics[scale=0.5]{AppendixA_Figures/AppendixA_d_NRC_Kstock_appendix_plot_Drates.png}
	\end{subfigure}
	\hfill
	\caption*{\tikz\draw[darkgreen, thick] (0,0) -- (0.5,0) node[right] {(dark green)}; Linear Time Trend}
	\hfill
	\caption*{\tikz\draw[blue, thick] (0,0) -- (0.5,0) node[right] {(blue)}; LOESS function Time Trend}
	\label{fig: d_ME_Kstock_appendix_plot_B}
\end{figure}
%In the case of the non-residential construction sector (${K^{g}}_{NRC}$), \citet{Tafunell2016}  does provide an index for its rate of growth in real terms, since 1900 year onwards. As previously mentioned,\citet{Tafunell2016} based on \citet{Tafunell2013} series and Ducoing (2012) doctoral dissertation, is built leveraging on data of cement inputs. While I asume this is not adequate for the construction sector as a whole, I do consider it a plausible assumption for the non-residential one. Thus, to sustain stock-flow consistency I estimate a new index for the growth of non-residential construction by using Tafunell et al. index as a proportion of the rate of growth of the gross fixed capital stock explained above, which allows also to estimate the residential sector as a reminder of the compound index implicit once considering Tafunell and Cilo-Lab PUC historical statistics. This index register an average and standard deviation statistical discrepancy with Tafunell et al (2015) of 1\% and 1.9\% respectively. To leverage on initial point estimates at $t=1900$, I assume that ${K^{n}}_{NRC}= \alpha {K^{n}}_{C}$  where $\alpha$  is the share of gross fixed capital formation in non-residential construction sector over the sum of residential and non-residential gross fixed capital formation available in Hofman (2000) table D, and follow the same assumption for the initial point estimate of total fixed capital consumption $D_{NRC} = \alpha D_{C}$.  From here, I got an initial point estimate at $t=1900$ for ${K^{g}}_{NRC}$ and then I use the compounded index constructed to sustain stock-flow consistency.  

For the data of gross fixed capital investment for the non-residential sector ${I^{g}}_{NRC}$ I leverage on gross fixed capital formation data in the construction sector as a whole from the consolidated date set by \citet{PerezEyzaguirre2019}, but preponderated by the share of investments implicit in Hofman (2000) series of gross fixed capital formation corresponding to the non-residential sector. In this case, the series do not rely on \citet{Tafunell2016} indexes to estimate gross fixed capital formation in the construction sector due the aforementioned assessment above. Therefore, I keep following a stock-flow consistency to build gross and net capital stocks with \citet{PerezEyzaguirre2019} extension of Clio-Lab Chilean Historical Statistics. 
	   
Retirements of capital stocks in the non-residential sector ($R_{NRC}$) that refers to fully depreciated value of plant infrastructure is estimated with stock-flow consistency by the equation $R_{NRC} = Ig_{NRC} - \Delta {K^{g}}_{NRC}$. Implicit annual depletion rate of non-residential construction ($z_{NRC}$) is determined with stock-flow consistency, and can be appreciated in the figure below. 


\begin{figure}[H]
	\centering
	\caption{Implicit Annual Depletion Rate of Gross Capital Stocks. Non-Residential Construction Sector (1901-1994).}
	\begin{subfigure}{0.5\textwidth}
		\centering
		\includegraphics[scale=0.5]{AppendixA_Figures/AppendixA_z_NRC_Kstock_appendix_plot_Zrates.png}
	\end{subfigure}
	\hfill
	\caption*{\tikz\draw[darkgreen, thick] (0,0) -- (0.5,0) node[right] {(dark green)}; Linear Time Trend}
	\hfill
	\caption*{\tikz\draw[blue, thick] (0,0) -- (0.5,0) node[right] {(blue)}; LOESS function Time Trend}
	\label{fig: d_ME_Kstock_appendix_plot_A}
\end{figure}

Departing from the initial point estimates for  ${K^{n}}_{NRC}$  ${D}_{NRC}$  with the assumptions presented above, and already having a series of gross capital stocks and retirements, I applied the identity coming out from the equations of motions previously mentioned and iteratively estimates of ${K^{n}}_{NRC}$ and also $D_{NRC}$. The implicit yearly depreciation rate ($d_{NRC}$) is estimated with the equation presented in the introduction of the appendix having as a result an average value of 2.16\%, a standard deviation of 0.73\%, and linear time trend slope of 0.86. This result seems consistent with the notion that depreciation of net NRC capital stocks should be characterized by higher variance in its flows of capital consumption than the construction sector as a whole, but lower than capital stocks in machinery and equipment. The figure presented below plots the time series of the depreciatation rate of net capital stocks for the non residential construction sector, suggesting an upward tendency that can be interpeted as an upward tendency of depreciation of the net value of plant infrastuctures, which is consistent with \cite{Harvey2018} understanding of built environments for the production process under a capital circuit framework. 

\begin{figure}[H]
	\centering
	\caption{Implicit Annual Depreciation Rate of Net Capital Stocks. Non-Residential Construction Sector (1901-1994).}
	\begin{subfigure}{0.5\textwidth}
		\centering
		\includegraphics[scale=0.5]{AppendixA_Figures/AppendixA_d_NRC_Kstock_appendix_plot_Drates.png}
	\end{subfigure}
	\hfill
	\caption*{\tikz\draw[darkgreen, thick] (0,0) -- (0.5,0) node[right] {(dark green)}; Linear Time Trend}
	\hfill
	\caption*{\tikz\draw[blue, thick] (0,0) -- (0.5,0) node[right] {(blue)}; LOESS function Time Trend}
	\label{fig: d_ME_Kstock_appendix_plot_B}
\end{figure}

\section{Measurements of Capital Stocks in Residential Construction Sector}

Given the compound index between Clio-Lab PUC data and \citet{Tafunell2016} index on gross fixed capital stock, provides a measure of the rate of growth of the non-residential construction sector as proportional rate of growth of the implicit index in Clio-Lab PUC series, all the variables for the residential construction sector are build as the reminder between the whole construction sector and the non-residential one. By doing so, the annual depreciation rates for residential construction are consistent with a similar patern of the construction sector as a whole and the non-residential construction sector in particular, of having a lower standard deviation at 2.86\% relative to its mean which registers at 5.04 \%, while presenting a high slope of its linear time trend at 5.62\% increase per year. Interestingly, this is consistent with \citet{Harvey2018} notion that the rate of turnover of built environment for social reproduction have a higher turnover than those for capitalist production ones. 

% --- Figure 1 ---
\begin{figure}[H]
    \centering
    \caption{\centering Implicit Annual Depletion Rate of Gross Capital Stocks, \\ Residential Construction Sector (1901–1994).}
    \includegraphics[width=0.5\textwidth]{AppendixA_Figures/AppendixA_z_RC_Kstock_appendix_plot_Zrates.png}
    \begin{minipage}{0.5\textwidth}
        \centering
        \footnotesize
        \tikz[baseline]{\draw[darkgreen, thick] (0,0)--(0.5,0);}~Linear trend (dark green) \hspace{1em}
        \tikz[baseline]{\draw[blue, thick] (0,0)--(0.5,0);}~LOESS trend (blue)
    \end{minipage}
    \label{fig:d_RC_Kstock_appendix_plot_A}
\end{figure}

% --- Figure 2 ---
\begin{figure}[H]
    \centering
    \caption{\centering Implicit Annual Depreciation Rate of Net Capital Stocks, \\ Residential Construction Sector (1901–1994).}
    \includegraphics[width=0.5\textwidth]{AppendixA_Figures/AppendixA_d_RC_Kstock_appendix_plot_Drates.png}
    \begin{minipage}{0.5\textwidth}
        \centering
        \footnotesize
        \tikz[baseline]{\draw[darkgreen, thick] (0,0)--(0.5,0);}~Linear trend (dark green) \hspace{1em}
        \tikz[baseline]{\draw[blue, thick] (0,0)--(0.5,0);}~LOESS trend (blue)
    \end{minipage}
    \label{fig:d_RC_Kstock_appendix_plot_B}
\end{figure}

%Given the compound index between Clio-Lab PUC data and \citet{Tafunell2016} index on gross fixed capital stock, provides a measure of the rate of growth of the non-residential construction sector as proportional rate of growth of the implicit index in Clio-Lab PUC series, all the variables for the residential construction sector are build as the reminder between the whole construction sector and the non-residential one. By doing so, the annual depreciation rates for residential construction are consistent with a similar patern of the construction sector as a whole and the non-residential construction sector in particular, of having a lower standard deviation at 2.86\% relative to its mean which registers at 5.04 \%, while presenting a high slope of its linear time trend at 5.62\% increase per year. Interestingly, this is consistent with Harvey's (2018) notion that the rate of turnover of built environment for social reproduction have a higher turnover than those for capitalist 

\begin{figure}[H]
	\centering
	\caption{Implicit Annual Depletion Rate of Gross Capital Stocks. Residential Construction Sector (1901-1994).}
	\begin{subfigure}{0.5\textwidth}
		\centering
		\includegraphics[scale=0.5]{AppendixA_Figures/AppendixA_z_RC_Kstock_appendix_plot_Zrates.png}
	\end{subfigure}
	\hfill
	\caption*{\tikz\draw[darkgreen, thick] (0,0) -- (0.5,0) node[right] {(dark green)}; Linear Time Trend}
	\hfill
	\caption*{\tikz\draw[blue, thick] (0,0) -- (0.5,0) node[right] {(blue)}; LOESS function Time Trend}
	\label{fig: d_RC_Kstock_appendix_plot_A}
\end{figure}


\begin{figure}[H]
	\centering
	\caption{Implicit Annual Depreciation Rate of Net Capital Stocks. Residential Construction Sector (1901-1994).}
	\begin{subfigure}{0.5\textwidth}
		\centering
		\includegraphics[scale=0.5]{AppendixA_Figures/AppendixA_d_RC_Kstock_appendix_plot_Drates.png}
	\end{subfigure}
	\hfill
	\caption*{\tikz\draw[darkgreen, thick] (0,0) -- (0.5,0) node[right] {(dark green)}; Linear Time Trend}
	\hfill
	\caption*{\tikz\draw[blue, thick] (0,0) -- (0.5,0) node[right] {(blue)}; LOESS function Time Trend}
	\label{fig: d_RC_Kstock_appendix_plot_B}
\end{figure}



\section{Measurements of Non-Residential and Total Capital Stocks}

For the case of non-residential capital stocks, and total capital stocks these are just constructed as the sum of machinery and equipment plus non-residential construction sector in the former, and the sum of every previous series constructed for the latter. Excersies of stock-flow consistency given the equation of motions and the identify for annual depreciation rates were implemented to check if additive properties and stock-flow consistency accounting matched in both cases to ensure the robustness of this time series successfully. Below can be appreciated the time series plot of the depreciation rate for both series, nevertheless at being implicit rates of adding different types of net capital stocks, their information is not immediately intuitive beyond that machinery and equipment is what should explain negative levels of depreciation. 

\vspace{-0.5cm}
\begin{figure}[H]
    \centering
    \caption{\centering Implicit Annual Depletion Rate of Gross Capital Stocks, \\ Non-Residential Sector (1901–1994).}
    \includegraphics[width=0.5\textwidth]{AppendixA_Figures/AppendixA_z_NR_Kstock_appendix_plot_Zrates.png}
     \begin{minipage}{0.5\textwidth}
        \centering
        \footnotesize
        \tikz[baseline]{\draw[darkgreen, thick] (0,0)--(0.5,0);}~Linear trend (dark green) \hspace{1em}
        \tikz[baseline]{\draw[blue, thick] (0,0)--(0.5,0);}~LOESS trend (blue)
    \end{minipage}
    \label{fig:d_NR_Kstock_appendix_plot}
\end{figure}

% --- Figure 1 ---
\begin{figure}[H]
    \centering
    \caption{\centering Implicit Annual Depletion Rate of Gross Capital Stocks, \\ Non-Residential Sector (1901–1994).}
    \includegraphics[width=0.5\textwidth]{AppendixA_Figures/AppendixA_z_NR_Kstock_appendix_plot_Zrates.png}
    \begin{minipage}{0.5\textwidth}
        \centering
        \footnotesize
        \tikz[baseline]{\draw[darkgreen, thick] (0,0)--(0.5,0);}~Linear trend (dark green) \hspace{1em}
        \tikz[baseline]{\draw[blue, thick] (0,0)--(0.5,0);}~LOESS trend (blue)
    \end{minipage}
    \label{fig:d_NR_Kstock_appendix_plot}
\end{figure}

% --- Figure 2 ---
\begin{figure}[H]
    \centering
    \caption{\centering Implicit Annual Depreciation Rate of Net Capital Stocks, \\ Non-Residential Sector (1901–1994).}
    \includegraphics[width=0.5\textwidth]{AppendixA_Figures/AppendixA_d_NR_Kstock_appendix_plot_Drates.png}
    \begin{minipage}{0.5\textwidth}
        \centering
        \footnotesize
        \tikz[baseline]{\draw[darkgreen, thick] (0,0)--(0.5,0);}~Linear trend (dark green) \hspace{1em}
        \tikz[baseline]{\draw[blue, thick] (0,0)--(0.5,0);}~LOESS trend (blue)
    \end{minipage}
    \label{fig:d_NR_Kstock_appendix_plot_A}
\end{figure}

% --- Figure 3 ---
\begin{figure}[H]
    \centering
    \caption{\centering Implicit Annual Depletion Rate of Gross Capital Stocks, \\ Total Sector (1901–1994).}
    \includegraphics[width=0.5\textwidth]{AppendixA_Figures/AppendixA_z_T_Kstock_appendix_plot_Zrates.png}
    \begin{minipage}{0.5\textwidth}
        \centering
        \footnotesize
        \tikz[baseline]{\draw[darkgreen, thick] (0,0)--(0.5,0);}~Linear trend (dark green) \hspace{1em}
        \tikz[baseline]{\draw[blue, thick] (0,0)--(0.5,0);}~LOESS trend (blue)
    \end{minipage}
    \label{fig:d_T_Kstock_appendix_plot_A}
\end{figure}

% --- Figure 4 ---
\begin{figure}[H]
    \centering
    \caption{\centering Implicit Annual Depreciation Rate of Net Capital Stocks, \\ Total Sector (1901–1994).}
    \includegraphics[width=0.5\textwidth]{AppendixA_Figures/AppendixA_d_T_Kstock_appendix_plot_Drates.png}
    \begin{minipage}{0.5\textwidth}
        \centering
        \footnotesize
        \tikz[baseline]{\draw[darkgreen, thick] (0,0)--(0.5,0);}~Linear trend (dark green) \hspace{1em}
        \tikz[baseline]{\draw[blue, thick] (0,0)--(0.5,0);}~LOESS trend (blue)
    \end{minipage}
    \label{fig:d_T_Kstock_appendix_plot_B}
\end{figure}









%%%%%%%%%%%%%%%%%%%%%% References %%%%%%%%%%%%%%%%%%%%%% 
\pagebreak
\bibliography{references.bib}